\documentclass[12pt,a4paper]{article}
\usepackage[fontset=fandol]{ctex}
\usepackage{amsmath,amssymb,amsthm}
\usepackage{geometry}
\usepackage{braket}

\geometry{left=2.5cm,right=2.5cm,top=2.5cm,bottom=2.5cm}

\title{表象变换}
\author{}
\date{}

\begin{document}
\maketitle

\section{离散态的表象变换}

设力学量 $\hat{A}$ 与 $\hat{B}$ 的本征方程分别为
\[
\hat{A}\ket{a_i}=a_i\ket{a_i},\qquad \hat{B}\ket{b_i}=b_i\ket{b_i},
\]
对应的本征矢各自满足完备性与正交归一条件：
\[
\sum_i\ket{a_i}\bra{a_i}=1,\qquad \sum_i\ket{b_i}\bra{b_i}=1,
\]
\[
\braket{a_i|a_j}=\delta_{ij},\qquad \braket{b_i|b_j}=\delta_{ij}.
\]

\subsection{变换算符}

定义表象变换算符
\[
\hat{S}=\sum_k\ket{a_k}\bra{b_k},
\]
则
\begin{align*}
\hat{S}\ket{b_j}&=\sum_k\ket{a_k}\braket{b_k|b_j}=\sum_k\ket{a_k}\delta_{kj}=\ket{a_j},\\[4pt]
\bra{a_j}\hat{S}&=\sum_k\braket{a_j|a_k}\bra{b_k}=\sum_k\delta_{jk}\bra{b_k}=\bra{b_j}.
\end{align*}

\subsection{态矢在不同表象中的表示}

任意态矢 $\ket{\psi}$ 在 $A$、$B$ 两表象中的列向量分别为
\[
\vec{\psi}_A=\begin{pmatrix}\braket{a_1|\psi}\\ \braket{a_2|\psi}\\ \vdots\end{pmatrix},
\qquad
\vec{\psi}_B=\begin{pmatrix}\braket{b_1|\psi}\\ \braket{b_2|\psi}\\ \vdots\end{pmatrix}.
\]

\subsection{变换矩阵的矩阵元}

在 $A$ 表象中，$\hat{S}$ 的矩阵元为
\begin{align*}
S_{ij}^{(A)}&=\bra{a_i}\hat{S}\ket{a_j}
=\sum_k\braket{a_i|a_k}\braket{b_k|a_j}
=\braket{b_i|a_j}.
\end{align*}

在 $B$ 表象中，$\hat{S}$ 的矩阵元为
\begin{align*}
S_{ij}^{(B)}&=\bra{b_i}\hat{S}\ket{b_j}
=\sum_k\braket{b_i|a_k}\braket{b_k|b_j}
=\braket{b_i|a_j}.
\end{align*}

\subsection{态矢的变换关系}

利用完备性条件，
\[
\braket{b_i|\psi}=\sum_k\braket{b_i|a_k}\braket{a_k|\psi}
=\sum_k S_{ik}\braket{a_k|\psi},
\]
即
\[
\boxed{\;\vec{\psi}_B=S\,\vec{\psi}_A\;}
\]

\subsection{力学量算符的变换关系}

算符 $\hat{H}$ 在 $B$ 表象中的矩阵元为
\begin{align*}
H_{ij}^{(B)}&=\bra{b_i}\hat{H}\ket{b_j}\\
&=\sum_{k,l}\braket{b_i|a_k}\bra{a_k}\hat{H}\ket{a_l}\braket{a_l|b_j}\\
&=\sum_{k,l}S_{ik}\,H_{kl}^{(A)}\,S_{lj}^{\dagger}.
\end{align*}
写成矩阵形式即
\[
\boxed{\;H^{(B)}=S\,H^{(A)}\,S^{\dagger}\;}
\]

\end{document}
